% ─────────────────────────────────────────────────
%  CHAPITRE 3 – CONCEPTION DE LA SOLUTION
% ─────────────────────────────────────────────────
\chapter{Conception de la solution}
\markboth{Conception de la solution}{Conception de la solution}

\section*{Introduction}
\phantomsection
\addcontentsline{toc}{section}{Introduction}

Ce chapitre traduit les besoins identifiés au chapitre précédent en une solution
structurée et implémentable. Il présente d'abord les principes retenus pour organiser
la plateforme, puis l'architecture logique d'I-SIB, les choix technologiques,
l'architecture technique, la conception détaillée du scénario de création d'un
prospect et enfin l'architecture de déploiement.

% ─────────────────────────────────────────────────
\section{Principes de conception de l'architecture}

La conception d'I-SIB repose sur la séparation des responsabilités et la
modularité. Chaque service prend en charge un domaine précis de la plateforme.

Cette organisation facilite la maintenance et l'ajout de nouvelles
fonctionnalités. Elle permet aussi de gérer les accès et de conserver une trace
des opérations importantes.

Ces choix conduisent à une architecture en microservices
\cite{fowler_microservices}, dans laquelle les services communiquent entre eux
selon les besoins.
% ─────────────────────────────────────────────────
\section{Architecture logique de la solution}


\subsection{Vue d'ensemble}

La solution I-SIB repose sur une architecture microservices organisée en plusieurs couches. Cette organisation permet de séparer les responsabilités entre les différents composants du système et de faciliter leur évolution.

L'architecture regroupe les interfaces utilisateur, les services métier et les composants techniques nécessaires au fonctionnement de la plateforme. La figure~\ref{fig:architecture_logique} présente l'organisation générale de ces éléments ainsi que leurs principales interactions.




\begin{figure}[H]
    \centering
    \MemoireFigureLarge{images/contenu/chapitre-3/architecture_fonctionnelle.png}
    \caption{Architecture logique de la plateforme I-SIB}
    \label{fig:architecture_logique}
\end{figure}

Comme le montre la figure~\ref{fig:architecture_logique}, la passerelle API joue
un rôle central : elle constitue la frontière entre les applications web et les
microservices internes. Les services métier restent ainsi isolés de tout accès
direct depuis l'extérieur.
\subsection{Description des couches et des blocs}

L'architecture logique est organisée en cinq couches. Chaque couche regroupe des
blocs qui ont un rôle proche. Cette organisation permet de séparer les
interfaces, la sécurité et le routage, les services métier, la gestion des
données ainsi que les outils nécessaires au fonctionnement de la plateforme.

\begin{description}
    \item[Couche présentation] Elle contient les interfaces utilisées pour accéder à
    la plateforme. Le \textbf{Front-Office} est l'interface destinée aux clients de
    la banque. Le \textbf{Back-Office} est utilisé par les employés.

Les administrateurs utilisent le \textbf{Module Manager} pour gérer les modules, les fonctionnalités disponibles et les droits associés. Les configurations définies depuis cette interface sont prises en charge par le \textbf{Backend Module Manager}, qui centralise les informations nécessaires au fonctionnement du système.

\item[Couche sécurité et routage] Cette couche contrôle l'accès à la plateforme et assure l'acheminement des requêtes vers les différents services. Le composant \textbf{Authentification et accès} gère l'identification des utilisateurs et la vérification de leurs autorisations. La \textbf{Passerelle de communication} constitue le point d'entrée des applications et redirige les requêtes vers les microservices concernés.

 Le \textbf{Registre des services} indique à la passerelle
    où se trouve chaque microservice disponible.

    \item[Couche métier] Elle regroupe les modules qui réalisent les principales
    fonctions de la plateforme : \textbf{CRM}, \textbf{SGI},
    \textbf{utilisateurs}, \textbf{comptes et transactions}, \textbf{agence},
    \textbf{comptabilité}, \textbf{trésorerie} et \textbf{crédit}. Le
    \textbf{service d'audit} enregistre les actions importantes.

    \item[Couche données et services techniques] La \textbf{base de données}
    conserve les informations, le \textbf{cache} accélère certains accès, la
    \textbf{messagerie} facilite les échanges entre services et le
    \textbf{stockage objet} contient les fichiers et les documents.

    \item[Infrastructure et exploitation] Cette couche rassemble les outils
    nécessaires au déploiement et au suivi de la plateforme. Les modules peuvent
    utiliser des \textbf{services d'intelligence artificielle} selon leurs
    besoins. Les outils d'\textbf{observabilité} servent à repérer les problèmes,
    tandis que la chaîne \textbf{CI/CD} automatise la construction et le
    déploiement des applications.
\end{description}

\section{Choix des outils et technologies}

Les technologies ont été choisies selon le rôle de chaque composant :
interfaces web, services métier, gestion des accès, échanges entre services,
stockage et déploiement. Cette section présente les principaux outils utilisés
dans I-SIB. Leur organisation sera ensuite détaillée dans l'architecture
technique.

Les interfaces web de la plateforme sont développées avec \textbf{Next.js 16}
\cite{nextjs_docs} et \textbf{TypeScript} \cite{typescript_docs}, afin de construire les applications front-office, back-office et Module Manager
modernes, maintenables et adaptées aux différents profils utilisateurs. Le backend repose
sur \textbf{Spring Boot 3.5} \cite{spring_boot_docs}, utilisé comme socle de développement des microservices
métier.

\textbf{Spring Cloud Gateway} \cite{spring_gateway_docs} sert de point d'entrée
aux applications. Il reçoit leurs requêtes et les transmet au microservice
concerné. \textbf{Keycloak} \cite{keycloak_docs} gère les utilisateurs et leurs
droits d'accès. \textbf{Netflix Eureka} \cite{spring_eureka_docs} permet à la
passerelle de retrouver les microservices disponibles.

Chaque microservice dispose de sa propre base \textbf{PostgreSQL}
\cite{postgresql_docs}. Cette organisation permet à chaque service de gérer ses données de manière indépendante. \textbf{Redis}
\cite{redis_docs} est utilisé pour certaines fonctions de cache et de limitation de débit. Les échanges asynchrones reposent sur \textbf{RabbitMQ}
\cite{rabbitmq_docs}, notamment pour la diffusion des événements d'audit. Les documents clients et les pièces KYC sont stockés dans \textbf{MinIO}
\cite{minio_docs}.

Les fonctionnalités d'intelligence artificielle s'appuient sur \textbf{Rasa}
\cite{rasa_docs} pour les interactions conversationnelles ainsi que sur des services dédiés à l'analyse et à l'évaluation des prospects.

 intervient pour les besoins de cache et certaines fonctionnalités techniques comme la limitation du nombre de requêtes. Les échanges asynchrones entre services reposent sur \textbf{RabbitMQ}
\cite{rabbitmq_docs}, notamment pour la diffusion des événements d'audit. Les documents clients et les pièces KYC sont quant à eux stockés dans \textbf{MinIO}
\cite{minio_docs}.

Les fonctionnalités d'intelligence artificielle reposent notamment sur \textbf{Rasa}
\cite{rasa_docs} pour les échanges conversationnels, ainsi que sur des services dédiés à l'évaluation des prospects et à l'analyse des données clients.

 La traçabilité technique est assurée par \textbf{Zipkin} \cite{zipkin_docs}, qui permet de
suivre les appels entre microservices. Enfin, \textbf{Docker} \cite{docker_docs} et
\textbf{Docker Compose} \cite{docker_compose_docs}
constituent l'environnement de référence pour le développement local et
l'organisation des composants en ensembles Compose.

\section{Architecture technique de la solution}


Après la présentation des technologies retenues, il est nécessaire de préciser leur organisation au sein de la plateforme.
L'architecture technique présente les principaux composants de la plateforme ainsi que leurs interactions.

Les applications front-office et back-office accèdent aux différents services par l'intermédiaire d'une passerelle API chargée de centraliser les accès. Les fonctionnalités métier sont réparties entre plusieurs microservices, chacun disposant de sa propre base de données. Des composants complémentaires, tels que la messagerie, le cache ou les outils de supervision, viennent compléter l'architecture.




Cette architecture permet de garantir la modularité du système et facilite son évolution. La figure~\ref{fig:architecture_technique} présente l’organisation générale des composants et leurs interactions.




\begin{figure}[H]
    \centering
    \MemoireFigureLarge{images/contenu/chapitre-3/architecture_technique.png}
    \caption{Architecture technique de la plateforme I-SIB}
    \label{fig:architecture_technique}
\end{figure}

% ─────────────────────────────────────────────────
\section{Conception détaillée}


Après la présentation des architectures logique et technique, cette section illustre le fonctionnement du module Relation Client à travers le scénario de création d’un prospect. Elle présente le périmètre retenu, le workflow associé ainsi que les principaux composants impliqués dans le traitement de la demande.
\subsection{Périmètre du scénario}


\subsection{Périmètre du scénario}

Le scénario étudié porte sur la création d'un prospect par un chargé de clientèle. Il comprend la saisie des informations, leur transmission au système, l'enregistrement du prospect dans le module CRM, l'attribution d'un score initial ainsi que l'enregistrement d'une trace d'audit. Les étapes suivantes du traitement du prospect ne sont pas prises en compte dans cette étude.

\subsection{Workflow de création}

Le processus commence lorsqu'un chargé de clientèle renseigne les informations du prospect depuis l'interface Back-office.

 La demande est ensuite transmise à la passerelle API, puis routée vers
le service CRM. Après validation des données, le prospect est enregistré, un score
initial est calculé par le service d’IA et une trace d’audit est générée. La
figure~\ref{fig:workflow_creation_prospect} présente les principales étapes de ce
processus.

\begin{figure}[H]
    \centering
    \MemoireFigureLarge{images/contenu/chapitre-3/workflow_creation_prospect.png}
    \caption{Workflow de création d'un prospect avec scoring IA}
    \label{fig:workflow_creation_prospect}
\end{figure}

{\footnotesize
\noindent
Les repères \textbf{1} à \textbf{3} couvrent la saisie, l'envoi vers la
passerelle API et le routage vers le service CRM. Le repère \textbf{4} indique
les contrôles métier. Les repères \textbf{5} à \textbf{8} montrent la
persistance, le calcul du score et la mise à jour du prospect. Le repère
\textbf{9} correspond à l'audit, tandis que les repères \textbf{10} à
\textbf{12} représentent le retour de la réponse jusqu'à l'affichage final. La
figure distingue ainsi le \textbf{Service CRM} de sa \textbf{Base de données CRM}.
\par}

\subsection{Synthèse des API}

Les opérations utiles au scénario sont exposées sous forme d'API REST sécurisées et
résumées dans le tableau~\ref{tab:endpoints_prospect}.

\begin{table}[H]
    \centering
    \caption{API mobilisées pour la création d'un prospect}
    \label{tab:endpoints_prospect}
    \footnotesize
    \renewcommand{\arraystretch}{1.08}
    \begin{tabularx}{\textwidth}{|p{1.7cm}|p{5.4cm}|X|}
        \hline
        \textbf{Méthode} & \textbf{Chemin} & \textbf{Utilisation principale} \\ \hline
        POST   & {\scriptsize\texttt{/api/crm/prospects}}                  & Enregistrer un nouveau prospect \\ \hline
        GET    & {\scriptsize\texttt{/api/crm/prospects/my-prospects}}     & Afficher le prospect créé dans la liste du chargé \\ \hline
        GET    & {\scriptsize\texttt{/api/crm/prospects/\{id\}}}           & Consulter la fiche créée et son score initial \\ \hline
    \end{tabularx}
\end{table}

\subsection{Vue de conception du scénario}

La figure~\ref{fig:conception_creation_prospect} complète le workflow en
reliant les interfaces \textit{boundary}, les composants de coordination
\textit{control} et les données métier \textit{entity}, selon les
stéréotypes de Jacobson \cite{jacobson_oose}.

\begin{figure}[H]
    \centering
    \MemoireFigureMedium{images/contenu/chapitre-3/conception_creation_prospect.png}
    \caption{Vue de conception du scénario Création d'un prospect}
    \label{fig:conception_creation_prospect}
\end{figure}

Ainsi, la conception détaillée de la création d'un prospect montre comment un
scénario métier s'appuie sur les couches définies précédemment : interface,
passerelle, service métier, données, intelligence artificielle et audit.

% ─────────────────────────────────────────────────
\section{Architecture de déploiement}

La plateforme est déployée sur un VPS AlmaLinux \cite{almalinux_official} à l’aide de Docker \cite{docker_docs} et Docker Compose \cite{docker_compose_docs}. Le déploiement est organisé en quatre ensembles : \texttt{infra}, \texttt{core}, \texttt{module} et \texttt{office}, qui regroupent respectivement les services communs, les référentiels, les modules métier et les interfaces utilisateur.

Deux réseaux assurent les communications : \texttt{sib\_public}, destiné aux accès externes, et \texttt{sib\_internal}, réservé aux échanges internes entre microservices et bases de données. La figure~\ref{fig:architecture_deploiement} présente cette organisation.
\begin{figure}[H]
    \centering
    \MemoireFigureLarge{images/contenu/chapitre-3/architecture_deploiement.png}
    \caption{Déploiement d'I-SIB sur un VPS AlmaLinux}
    \label{fig:architecture_deploiement}
\end{figure}

Dans la pratique, le déploiement est organisé par ensembles. L'ensemble
\texttt{infra} regroupe les services de base nécessaires au fonctionnement de
la plateforme, comme la passerelle API, Keycloak, Eureka,
Redis, RabbitMQ, etc. Les ensembles \texttt{core} et \texttt{module} regroupent les
services applicatifs, tandis que \texttt{office} contient les interfaces. Lorsqu'une
modification est intégrée au projet, GitHub Actions \cite{github_actions_docs}
déclenche le redéploiement des stacks concernées via Portainer
\cite{portainer_docs}, ce qui facilite le suivi du déploiement sur le VPS.
% ─────────────────────────────────────────────────
\section*{Conclusion}
\phantomsection
\addcontentsline{toc}{section}{Conclusion}


Ce chapitre a présenté la conception fonctionnelle et technique d'I-SIB ainsi que le scénario de création d'un prospect. Il a également décrit le déploiement de la solution sur le VPS AlmaLinux \cite{almalinux_official}. Le chapitre suivant présente les résultats obtenus.


